\section{Introduction}
Among the three-dimensional model geometries, $\Solv$, i.e. $\mathbb{R}^3$ equipped with the metric $e^{2z}dx^2+e^{-2z}dy^2+dz^2$, is the least symmetric one~\cite{Sco}. This makes $\Solv$-manifolds (i.e. compact $3$-manifolds admitting a $\Solv$ metric) a very special class within the classification scheme of Thurston's geometrization theorem; in fact, they can be characterized as the geometric manifolds which are neither Seifert nor hyperbolic. From a historical perspective, their importance stems from the fact that many $\Solv$ manifolds arise as cusps of Hirzebruch modular surfaces~\cite{Hir}; and the understanding of their signature defect was the main motivation behind the discovery of the Atiyah--Patodi--Singer index theorem for manifolds with boundary~\cite{APS}, see~\cite{ADS}. In a related fashion, three-dimensional $\Solv$ manifolds are also among the simplest examples where non-abelian Fourier analysis can be performed~\cite{Bre}. More recently, the computation of their Heegaard Floer homology has provided evidence for the far-reaching $L$-space conjecture~\cite{BCW}.


In this paper we study the monopole Floer homology of a $\Solv$ rational homology sphere $Y$ from a geometric viewpoint. Monopole Floer homology is a package of invariants of three-manifolds introduced by Kronheimer and Mrowka in~\cite{KM} obtained by studying the Seiberg--Witten equations (see also~\cite{Lin3} for a friendly introduction). While monopole Floer homology is a topological invariant, and can be therefore computed in many cases using tools such as surgery exact triangles~\cite{KMOS}, it is interesting to understand its relation with special geometric structures on the space, the case of Seifert fibered spaces~\cite{MOY} being the prototypical example. In our case, a $\Solv$-rational homology sphere $Y$ has the structure of a torus semibundle, and admits several different $\Solv$-metrics obtained by rescaling the metrics along the fibers (see Section~\ref{fourier} for a more detailed discussion of $\Solv$ geometry). Our main result is then the following.
\begin{theo}\label{main}
Let $Y$ be a $\Solv$-rational homology sphere, equipped with a $\Solv$ metric. If the fibers are small enough, then there are small regular perturbations for which the Seiberg--Witten equations on $Y$ do not admit irreducible solutions.
\end{theo}

The following is an immediate consequence of the Theorem~\ref{main}. Recall that a rational homology sphere $Y$ is an $L$-\emph{space} if $\HMf_*(Y,\spin)=\mathbb{Z}[U]$ as a $\mathbb{Z}[U]$-module for each spin$^c$ structure $\spin$.


\begin{coro}\label{corol}
Let $Y$ be a $\Solv$-rational homology sphere. Then $Y$ is an $L$-space.
\end{coro}

The analogous result in the setting of Heegaard Floer homology (which is known to yield isomorphic invariants, see~\cite{CGH1,HFHM1} and subsequent papers) was proved by topological means in~\cite{BCW} with $\mathbb{Z}/2\mathbb{Z}$-coefficients, and extended to $\mathbb{Z}$-coefficients in~\cite{RR}. Let us also point out that compact $\Solv$ manifolds have either $b_1=0$ or $1$; in the latter case, they are Anosov torus bundles over the circle, and their Heegaard Floer homology (with $\mathbb{Z}$ coefficients) was computed in~\cite{Bal}.


In our approach, we look at the monopole Floer homology of $\Solv$-manifolds from the point of view of spectral geometry. The main ingredient in the proof of Theorem~\ref{main} is the following relation, for a rational homology sphere, between the existence of irreducible solutions to the Seiberg--Witten equations and the first eigenvalue $\lambda_1^*$ of the Hodge Laplacian on coexact $1$-forms (which improves on the main result of~\cite{Lin4}).


\begin{theo}[Theorem~$3$ of~\cite{LL}]\label{spectral}
Let $Y$ be a rational homology sphere equipped with a metric $g$. Denote by $\tilde{s}(p)$ the sum of the two least eigenvalues of the Ricci curvature at the point $p$. If the inequality
\[
\lambda_1^*\geq -\inf_{p\in Y}\tilde{s}\frac{(p)}{2}
\]
holds, then the Seiberg--Witten equations do not admit irreducible solutions.
\end{theo}

In the case of a $\Solv$-metric, $\tilde{s}=-2$ at every point, so in order to prove Theorem~\ref{main}, we need to show that for suitable $\Solv$-metrics on $Y$, $\lambda_1^*\geq 1$. Let us describe the strategy behind the proof of this by discussing the content of each section.


In Section~\ref{fourier}, we review some facts about the geometry and topology of $\Solv$-manifolds. As $\Solv$ is the left-invariant metric for a solvable Lie group structure on $\mathbb{R}^3$, one can study Fourier analysis on it, and we will introduce the basic ideas behind it. In Section~\ref{spectralbound}, we use the aforementioned Fourier analysis to show that, for metrics with sufficiently small fibers, $\lambda_1^*= 1$, so that the Seiberg--Witten equations do not admit irreducible solutions by Theorem~\ref{spectral}. As these metrics have $\lambda^*_1$ is exactly $1$, they lie in the borderline case of Theorem~\ref{spectral}, and transversality is a quite subtle issue. We discuss it in Section~\ref{trans}, where we will study explicit small perturbations of the equations and existence of harmonic spinors.



%\vspace{0.3cm}
\section{Compact Solvmanifolds and their Fourier analysis}\label{fourier}

We start by reviewing the basics of $\Solv$-geometry; most of the following discussion is taken from Section~$12.7$ of~\cite{Mar}. Recall that $\Solv$ is the Riemannian manifold $\mathbb{R}^3$ equipped with the metric
\[
e^{2z}dx^2+e^{-2z}dy^2+dz^2.
\]
This is the left-invariant Riemannian metric on $\mathbb{R}^3$ when equipped with the solvable Lie group structure
\[
(x,y,z)\cdot(x',y',z')=(x+e^{-z}x', y+e^{z}y', z+z').
\]
This can be though of as the semidirect product corresponding to the splitting of
\[
0\rightarrow \mathbb{R}^2\rightarrow \Solv\stackrel{p}{\longrightarrow} \mathbb{R}\rightarrow 0,
\]
where $p(x,y,z)=z$, given by
\[
z\mapsto
\begin{bmatrix}e^{-z}&0\\0 &e^{z}
\end{bmatrix}\in\SL(2,\mathbb{R}),
\]
seen as linear automorphisms of $\mathbb{R}^2$. The Ricci tensor is given in this coordinates by
\[
\begin{bmatrix}0&0&0\\0 &0&0\\ 0&0&-2
\end{bmatrix},
\]
so that both $s$ and $\tilde{s}$ are $-2$ at each point. We can see that the foliation in $\mathbb{R}^2$ by the planes with $z$ constant descend to any compact $\Solv$-manifold; in fact, it descends to a foliation for which all the leaves are tori or Klein bottles.
%\\


Orientable compact solvmanifolds either have $b_1=0$ or $1$. The manifolds of the latter type, which will be denoted by $\tilde{Y}$, arise as quotients $\Gamma\setminus\Solv$ for lattices $\Gamma\subset \Solv$. Every such lattice is a split extension
\[
0\rightarrow \Lambda\rightarrow \Gamma\stackrel{p}{\longrightarrow} a\mathbb{Z}\rightarrow 0,
\]
where $\Lambda\subset \mathbb{R}^2$ is a lattice invariant under the action of $
\begin{bmatrix}e^{-a}&0\\0 &e^{a}
\end{bmatrix}$. The underlying topological manifold is a torus bundle with monodromy $A\in \SL(2,\mathbb{Z})$; here $|\tr A|>2$ (i.e. $A$ is Anosov) and $e^a$ and $e^{-a}$ are its eigenvalues.

\begin{exam}\label{figure8}
Consider $ A=
\begin{bmatrix}2&1\\1 &1
\end{bmatrix}.$

The mapping torus is well-known to be the zero surgery on the figure eight knot. Its eigenvalues are $\varphi^2$ and $\varphi^{-2}$ where $\varphi=\frac{1+\sqrt{5}}{2}$ is the golden ratio. Recall that it satisfies $\varphi^2=\varphi+1$. Consider the vectors
\[
v=(1-\varphi,\varphi)\quad w=(1,1).
\]
If $S$ is the matrix with columns $v$ and $w$, we have
\[
A=S^{-1}
\begin{bmatrix}\varphi^{-2}&0\\0 &\varphi^{2}
\end{bmatrix}S;
\]
setting $\Lambda$ to be the lattice generated by $v$ and $w$, and $a=\log(\varphi^2)$, we obtain the lattice $\Gamma$ equipping the mapping torus of $A$ with a $\Solv$ metric.
\end{exam}

\begin{rema}\label{NT}
We can also think about this construction from a more number theoretic viewpoint, which makes the connection with~\cite{Hir} and~\cite{ADS} clearer. Consider the field $k=\mathbb{Q}(\sqrt{5})$. It is totally real, and it comes with two natural embeddings $\phi_+,\phi_-$ into $\mathbb{R}$ sending $\sqrt{5}$ to $\pm\sqrt{5}$. The ring of integers $\mathcal{O}_k$ is the lattice $\Lambda=\mathbb{Z}[\varphi]$ which has basis $\varphi$ and $1$. The group of totally positive units is generated by $\varphi^2$; and it is easy to see that its multiplication action is given in our chosen basis by $A$. Finally, we can embed the lattice $\Lambda$ in $\mathbb{R}^2$ using $(\phi_-,\phi_+)$; our basis elements are mapped to the vectors $v$ and $w$.


This construction is readily generalized to any monodromy $A\in \SL(2,\mathbb{Z})$ by looking at the action of a totally positive or negative unit on fractional ideals in real quadratic number fields.
\end{rema}

A $\Solv$-manifold with $b_1=0$, denoted by $Y$, is a torus semibundle; therefore it admits a double cover $\tilde{Y}$ which is a $\Solv$ torus bundle $\Gamma\setminus\Solv$. Denoting the cover involution on $\tilde{Y}$ by $\tau$, we can describe $Y$ in the following way. There is a basis $v,w$ of the lattice $\Lambda=\Gamma\cap \mathbb{R}^2$, for which $\tau$ is the order $2$ isometry
\[
(av+bw,z)\mapsto \left(\left(a+\frac{1}{2}\right)v-bw,-z\right).
\]
Recall that there are only two orientation preserving isometries of $\Solv$ fixing the origin and inducing an orientation reversing isometry on $\mathbb{R}^2\times \{0\}$, namely
\[
(x,y,z)\mapsto (\pm y, \pm x, -z).
\]
In particular, the fact that $\tau$ is an isometry implies that $v$ and $w$ are multiples of $(1,1)$ and $(1,-1)$, or viceversa.

\begin{rema}\label{mod2}
The existence of such an involution on $Y$ provides non trivial constraints on the monodromy $A$. Among the others, we have that $A$ is congruent to the identity modulo $2$. This readily follows from the fact that the action is fixed point free (see~\cite[Section~6]{Sak}).
\end{rema}

From this description, we see that on any $\Solv$-manifold we obtain a one parameter family of metrics obtained by rescaling the lattice $\Lambda$; this can be seen concretely in Example~\ref{figure8}.
%\\


Let us now introduce the basics of Fourier analysis on a compact Solvmanifold with $b_1(Y)=1$. We follow~\cite[the first chapter]{Bre}, to which we refer for a pleasant, more thorough, discussion.


Consider a smooth function $f:\Gamma\setminus\Solv\rightarrow\mathbb{R}$. This can be thought (with a little abuse of notation) as a function $f:\Solv\rightarrow \mathbb{R}$ which is left invariant under $\Gamma$. In particular, it is invariant under the action of $\Lambda\subset\Gamma$, i.e.
\[
f(\underline{x}+m,z)=f(\underline{x},z)\text{ for all }m\in \Lambda.
\]
We can therefore expand $f$ in Fourier series in the $\mathbb{R}^2\times\{0\}\subset\Solv$ directions
\[
f(\underline{x},z)=\sum_{\mu\in \Lambda'} a_{\underline{\mu}}(z)e^{ i \underline{\mu}\cdot \underline{x}}.
\]
for some smooth functions $a_{\underline{\mu}}(z)$. Here $\Lambda'$ is the dual lattice of $\Lambda$, where we use the convention
\[
\Lambda'=\left\{\underline{\mu}\in \mathbb{R}^2\middle| \underline{\mu}\cdot m\in 2\pi\mathbb{Z}\text{ for all }m\in\Lambda\right\}.
\]
We now use the fact that $f$ is invariant by the action of $(\underline{0},a)$. Letting $A=
\begin{bmatrix}e^{-a}&0\\0 &e^{a}
\end{bmatrix}$, we see that
\[
f(\underline{x},z)=f((\underline{0},a)\cdot(\underline{x},z))=f(A\underline{x}, z+a),
\]
hence, after reindexing,
\[
\sum_{\underline{\mu}\in \Lambda'} a_{\underline{\mu}}(z)e^{ i \underline{\mu}\cdot \underline{x}}=\sum_{\underline{\mu}\in \Lambda'} a_{\underline{\mu}}(z+a)e^{ i \underline{\mu}\cdot A\underline{x}}=\sum_{\underline{\mu}\in \Lambda'} a_{\underline{\mu}\cdot A}(z+a)e^{ i \underline{\mu}\cdot\underline{x}}.
\]
This implies that
\[
a_{\underline{\mu}}(z)=a_{\underline{\mu}\cdot A}(z+a),
\]
so $a_{\underline{\mu}}$ determines via translation $a_{\underline{\mu}\cdot A^n}$. In particular, the Fourier series is determined by the collection of functions for $a_{\underline{\mu}}(z)$ for $\underline{\mu}\in \Lambda'/V$, $V$ being the group of automorphisms of the dual lattice $\Lambda'$ generated by $A$. While $a_0$ is a periodic function with period $a$, it can be shown that the functions $a_{\underline{\mu}}(z)$ for $\underline{\mu}\neq 0$ are in the Schwartz--type space
\begin{equation}\label{schwartz}
\mathcal{S}=\left\{f\middle| e^{nz}f^{(m)}(z) \text{ is bounded for all }n\in\mathbb{Z},m\geq0\right\},
\end{equation}
where $f^{(m)}$ denotes the $m^{\rm th}$ derivative of $f$.
%\\


With this in mind, let us study as a warm-up example the Laplacian on functions on $\Gamma\setminus\Solv$, which can be written as
\[
\Delta f=-\left(e^{-2z}f_{xx}+e^{2z}f_{yy}+f_{zz}\right).
\]
Let us use the decomposition in Fourier modes discussed above. We then have a $L^2$-unitary decomposition
\[
\Delta=\bigoplus_{\underline{\mu}\in \frac{\Lambda'}{V}} \Delta_{\underline{\mu}},
\]
where $\Delta_{0}$ acts on $L^2(\mathbb{R}/a\mathbb{Z})$ and $\Delta_{\underline{\mu}}$ is a diagonalizable operator on $L^2(\mathbb{R})$. In particular, if we have $\underline{\mu}=(\mu,\mu')$, the corresponding operator is given by substituting
\[
\frac{d}{dx}\mapsto i\mu,\quad\frac{d}{dy}\mapsto i\mu'
\]
so that
\[
\Delta_{\underline{\mu}} f=-f_{zz}+\left(\mu^2 e^{-2z}+(\mu')^2e^{2z}\right)f.
\]
Therefore $\lambda^2$ is an eigenvalue of $\Delta_{\underline{\mu}}$ if and only if
\[
f_{zz}=\left(\mu^2 e^{-2z}+(\mu')^2e^{2z}-\lambda^2\right)f
\]
admits a non-zero solution in $L^2(\mathbb{R})$. While this equation is not solvable in terms of elementary functions, we can still understand the basic properties of its spectrum. Let us first recall the following well-known elementary Lemma~\ref{keyODE}.

\begin{lemm}\label{keyODE}
Suppose $f:\mathbb{R}\rightarrow\mathbb{R}$ solves the second order linear ODE
\[
f_{zz}=\Phi(z)\cdot f
\]
where $\Phi$ is smooth and $\Phi(z)>0$ everywhere. If $f$ is not identically zero, $f$ cannot be in $L^2(\mathbb{R})$.
\end{lemm}

\begin{proof}
Possibly after replacing $f(z)$ by $-f(z)$ or $f(-z)$, we can assume that at $x_0$ both $f(x_0)=c>0$ and $f'(x_0)\geq0$. Suppose there is $t_0>x_0$ with $0<f(t_0)<f(x_0)$. We can also assume $f>0$ on $[x_0,t_0]$. Then there is $x_0<t<t_0$ with $f'(t)<0$. Applying again the mean value theorem, there is $x_0<t'<t$ with $f''(t')<0$, which is contradiction as $f''(t')=\Phi(z)\,\cdot\,f>0$. So $f(x)\geq c$ for $x\geq x_0$, and the result follows.
\end{proof}

We then have the following.

\begin{lemm}\label{eigenvalue}
For $\underline{\mu}\neq0$ the first eigenvalue of $\Delta_{\underline{\mu}}$ is at least $2|\mu\mu'|\neq0$.
\end{lemm}


\begin{proof}
By \textsc{AM-GM}, the inequality
\[
\mu^2 e^{-2z}+(\mu')^2e^{2z}\geq 2|\mu\mu'|,
\]
holds, and the result follows from the previous Lemma~\ref{eigenvalue}.
\end{proof}

In terms of the number theoretic description in Remark~\ref{NT}, the quantity $\mu\mu'$ is the \emph{norm} $N(\underline{\mu})$. The only basic property we will need is that there is $c>0$ such that $|\mu\mu'|\geq c$ for all $\underline{\mu}\in \Lambda'\setminus \{0\}$; for example, we can choose $c=1$ in Remark~\ref{NT}.
%\\


For completeness, let us conclude this section by discussing the zero mode $\underline{\mu}=0$. In this case, we study the ODE
\[
f_{zz}=-\lambda^2 f
\]
with $f$ periodic with period $a$. It has eigenvalues $\lambda^2=\frac{4\pi^2}{a^2} n^2$ for $n\in\mathbb{Z}$.


%\vspace{0.3cm}

\section{The spectrum on coexact \texorpdfstring{$1$}{1}-forms}\label{spectralbound}
In this section we will perform the key computation behind our main result. Recall from the previous section that on a $\Solv$-manifold there is a non-trivial family of metrics obtained by rescaling the lattice $\Lambda\subset \mathbb{R}^2$. With this in mind, we have the following.

\begin{prop}\label{maineigen}
Let $Y$ be a rational homology sphere equipped with a $\Solv$ metric such that the fibers are small enough. Then the first eigenvalue of $\Delta$ on coexact $1$-forms satisfies $\lambda_1^*=1$. Furthermore, the $1$-eigenspace is one dimensional.
\end{prop}

In fact, our proof will provide an explicit smallness condition for the fibers.
%\\


Let us start by considering the case of a $\Solv$-manifold $\tilde{Y}=\Gamma\setminus \Solv$ with $b_1=1$. The $1$-forms
\[
\X=e^zdx,\quad \Y=e^{-z}dy\quad \Z=dz
\]
descend to a left-invariant dual orthonormal frame on $\tilde{Y}$. We can then write any $1$-form $\xi$ as
\[
\xi=f\X+g\Y+h\Z,
\]
where $f,g,h$ are functions on $\Gamma\setminus\Solv$, or equivalently left-invariant functions on $\Solv$. We are interested in understanding for which $\lambda\in\mathbb{R}$ the equation
\[
\ast d\xi=\lambda\xi
\]
admits non-trivial solutions. Notice that, provided $\lambda\neq0$, such a form necessarily satisfies $d\ast\xi=0$, i.e. it is coclosed. In particular, $\lambda^2$ is an eigenvalue of the Laplacian on coexact $1$-forms. We have
\[
d\xi=\left(e^{-z}g_x-e^zf_y\right)\X\wedge \Y+
\left(-g_z+g+e^zh_y\right)\Y\wedge \Z+
\left(f_z+f-e^{-z}h_x\right)\Z\wedge \X
\]
so that our equation is equivalent to the system
\begin{equation}\label{eqn}
\begin{split}
-g_z+g+e^zh_y&=\lambda f\\
f_z+f-e^{-z}h_x&=\lambda g\\
e^{-z}g_x-e^zf_y&=\lambda h, 
\end{split}
\end{equation}
while coclosedness is equivalent to
\[
e^{-z}f_x+e^zg_y+h_z=0.
\]
Differentiating we get
\begin{align*}
-e^{-2z}h_{xx}&=-e^{-z}f_{xz}-e^{-z}f_x+\lambda e^{-z}g_x\\
-e^{2z}h_{yy}&=-e^{z}g_{yz}+e^z g_y-\lambda e^{z}f_y\\
-h_{zz}&=e^{-z}f_{xz}-e^{-z}f_x+e^zg_{yz}+e^zg_y,
\end{align*}
therefore summing we obtain
\[
\Delta h=\lambda^2 h-2e^{-z}f_x+2e^zg_y,
\]
where $\Delta$ denotes the Laplacian on functions on $\tilde{Y}$. Similarly for $g$ we obtain
\begin{align*}
-e^{-2z}g_{xx}&=-\lambda e^{-z} h_x-f_{xy}\\
-e^{2z}g_{yy}&=f_{xy}+e^zh_{yz}\\
-g_{zz}&=\lambda f_z-g_z-e^zh_y-e^zh_{yz},
\end{align*}
hence summing
\begin{align*}
\Delta g&=\lambda(f_z-e^{-z}h_x)-g_z-e^zh_y=\lambda^2g-\lambda f-g_z-e^zh_y=\\
&=\left(\lambda^2-1\right)g-2e^z h_y.
\end{align*}
Finally, as
\begin{align*}
-e^{-2z}f_{xx}&=e^{-z}h_{xz}+g_{xy}\\
-e^{2z}f_{yy}&=\lambda e^z h_y-g_{xy}\\
-f_{zz}&=f_z-e^{-z}h_{xz}+e^{-z}h_x-\lambda g_z,
\end{align*}
we have
\begin{align*}
\Delta f&=\lambda(-g_z+e^zh_y)+f_z+e^{-z}h_x=\lambda^2 f-\lambda g+f_z+e^{-z}h_x=\\
&=\left(\lambda^2-1\right)f+2e^{-z}h_x.
\end{align*}
Notice that $\Z$ is a harmonic $1$-form; as $b_1=1$, all harmonic forms are multiples of~it.

\begin{lemm}\label{cover}
Let $\tilde{Y}$ be a $\Solv$ manifold with $b_1=1$ equipped with a metric for which the fibers are small enough. Then $\lambda_1^*=1$, and the $1$-eigenspace is spanned by $\mathcal{X}$ and $\mathcal{Y}$.
\end{lemm}


\begin{proof}
We can expand $f,g$ and $h$ in Fourier series; the operator $\ast d$ decomposes accordingly in the sum of $\ast d_{\underline{\mu}}$, and in the $\underline{\mu}$ component our equations look like
\begin{align*}
\Delta_{\underline{\mu}}h&=\lambda^2 h-2i\mu e^{-z} f+2i\mu'e^z g\\
\Delta_{\underline{\mu}}g&=\left(\lambda^2-1\right)g-2i\mu'e^zh\\
\Delta_{\underline{\mu}}f&=\left(\lambda^2-1\right)f+2i\mu e^{-z}h
\end{align*}
with $f,g$ and $h$ are \emph{complex} valued functions in the space $\mathcal{S}$.


Let us discuss first the modes $\underline{\mu}\neq0$. By Lemma~\ref{eigenvalue}, the bottom of the spectrum of $\Delta_{\underline{\mu}}$ is bounded below by $2|\mu\mu'|$; and furthermore, by suitably rescaling the metric, we can arrange that this quantity is $>16$ for all $\underline{\mu}\neq0$. Multiplying each equation by $\bar{h},\bar{g}$ and $\bar{f}$ respectively, and adding them together, we obtain the pointwise identity
\begin{multline*}
\bar{h}\Delta_{\underline{\mu}}h+\bar{g}\Delta_{\underline{\mu}}g+\bar{f}\Delta_{\underline{\mu}}f\\
= \lambda^2|h|^2+\left(\lambda^2-1\right)|g|^2+\left(\lambda^2-1\right)|f|^2+4\Rel\left(i\mu'e^zg\bar{h}\right)-4\Rel\left(2i\mu e^{-z}f\bar{h}\right).
\end{multline*}
In particular, this implies that the left-hand side is real. By the Peter--Paul inequality, we have the pointwise inequalities
\begin{align*}
\left|4\Rel(i\mu'e^zg\bar{h})\right|&\leq4\left|\mu'e^z\bar{h}\right|\,|g|\leq \frac{(\mu')^2e^{2z}}{2}|h|^2+8|g|^2\\
\left|4\Rel(i\mu e^{-z}f\bar{h})\right|&\leq4\left|\mu e^{-z}\bar{h}\right|\,|f|\leq\frac{\mu^2e^{-2z}}{2}|h|^2+8|f|^2
\end{align*}
so that
\begin{equation}\label{ineq}
\bar{h}\tilde{\Delta}_{\underline{\mu}}h+\bar{g}\Delta_{\underline{\mu}}g+\bar{f}\Delta_{\underline{\mu}}f\leq \lambda^2|h|^2+(\lambda^2+7)|g|^2+(\lambda^2+7)|f|^2
\end{equation}
where
\[
\tilde{\Delta}_{\underline{\mu}}h=-h_{zz}+\frac{1}{2}\left(\mu^2 e^{-2z}+(\mu')^2e^{2z}\right)h
\]
is still a diagonalizable operator over $L^2(\mathbb{R})$. The same argument as Lemma~\ref{eigenvalue} implies that the first eigenvalue of $\tilde{\Delta}_{\underline{\mu}}$ is at least $|\mu\mu'|$. Therefore, by integrating the Inequality~\eqref{ineq} we have
\[
|\mu\mu'|\left(\|h\|^2+\|g\|^2+\|f\|^2\right)\leq \left(\lambda^2+7\right)\left(\|h\|^2+\|g\|^2+\|f\|^2\right).
\]
As by assumption $|\mu\mu'|>8$, $\lambda^2>1$.


Finally, we deal with the zero mode. Suppose $0<\lambda^2< 1$. Then $\lambda^2-1<0$, hence
\[
-g_{zz}=\left(\lambda^2-1\right)g,\quad -f_{zz}=\left(\lambda^2-1\right)f
\]
have no non-zero periodic solution. It follows from Equation~\eqref{eqn} that $h$ is constant, so we have a multiple of the harmonic form $\Z$. Finally, the case $\lambda^2=1$ corresponds to the span of $\X$ and $\Y$.
\end{proof}

Finally, we are ready to prove Proposition~\ref{maineigen}.
\begin{proof}[Proof of Proposition~\ref{maineigen}]
Suppose $Y$ is a $\Solv$-rational homology sphere. Consider its double cover $\pi:\tilde{Y}\rightarrow Y$ where $\tilde{Y}$ has $b_1(\tilde{Y})=1$. If $\xi$ is a $\lambda$-eigenform on $Y$, the $\pi^*\xi$ is a $\lambda$-eigenform on $\tilde{Y}$. Choose a $\Solv$-metric with fibers small enough, so that Lemma~\ref{cover} applies. This implies that on $Y$ we have $\lambda_1^*\geq1$, and furthermore that if $\xi$ is a $1$-eigenform on $Y$, then $\pi^*\xi$ is a linear combination of $\X$ and $\Y$. Finally, in the notation of Section~\ref{fourier}, if $v,w$ is the basis of $\Lambda$, then exactly the linear combinations of $\X$ and $\Y$ that vanish on $w$ at $z=0$ descend to $Y$.
\end{proof}


\begin{rema}
As the covering involution $\tau$ of $\tilde{Y}$ sends $\X$ to $\pm \Y$, the forms that descend are multiples of either $\X+\Y$ or $\X-\Y$.
\end{rema}


%\vspace{0.3cm}
\section{Transversality}\label{trans}
In the previous section, we have exhibited a metric for which $\lambda_1^*=-\inf (\tilde{s}/2)$. As this is the borderline case of Theorem~\ref{spectral}, transversality is a quite delicate issue as small perturbation might introduce irreducible solutions. This should be compared with the discussion of flat manifolds in~\cite[Chapter~37]{KM}. As in their setting, we will show that we can achieve transversality, while still not having irreducible solutions, by considering the perturbed functional
\[
\mathcal{L}(B,\Psi)-\frac{\delta}{2}\|\Psi\|^2
\]
for $\delta$ sufficiently small. The corresponding equations for the critical points are
\begin{align*}
D_B\Psi&=\delta\Psi\\
\frac{1}{2}\rho(F_{B^t})&=(\Psi\Psi^*)_0.
\end{align*}
We will denote by $\eta$ the unique unit length $1$-eigenform such that $\eta(v)>0$ at $z=0$ and $\eta$ descends to $Y$. Recall (\cite[Lemma~1.4]{Tur}) that there is a natural one-to-one correspondence between spin$^c$ structures and unit length $1$-forms up to homotopy outside balls; denote by $\spin_0$ the spin$^c$ structure on $Y$ corresponding to $\eta$. With this in mind, we have the following.


\begin{lemm}
Consider a spin$^c$ structure $\spin\neq \spin_0, \bar{\spin}_0$. Then, for $\delta$ small enough, the perturbed Seiberg--Witten equations do not admit irreducible solutions.
\end{lemm}

\begin{proof}
Suppose we have a sequence $\delta_i\rightarrow 0$ with corresponding irreducible solutions $(B_i,\Psi_i)$; consider the corresponding configurations in the blow-up $(B_i,r_i,\psi_i)$, where $\|\psi_i\|_{L^2}=1$. These admit (up to gauge transformations, and up to passing to a subsequence) a limit $(B,r,\psi)$ which solves the blown-up equations with $\delta=0$; in particular, as the unperturbed equations do not admit irreducible solutions by Theorem~\ref{spectral}, $r=0$, $B$ is the flat connection, and $D_B\psi=0$. Recall that, setting $\xi=\rho^{-1}(\Psi\Psi^*)_0$, it is shown in~\cite[Proposition~$1$]{LL} that for solutions $(B,\Psi)$ of the \emph{unperturbed} Seiberg--Witten equations the pointwise identity
\[
|\nabla\xi|^2+|d\xi|^2=|\Psi|^2|\nabla_B\Psi|^2
\]
holds. This holds for the perturbed equations up to an error going to zero for $\delta_i\rightarrow 0$; hence it will apply to the limit form $\alpha=\rho^{-1}(\psi\psi^*)_0$. Furthermore, as it is the limit of the sequence of coexact forms $\frac{1}{2r_i^2}\ast F_{B^t}$, $\alpha$ is a coexact $1$-form.


Let us study the geometry of $\alpha$. As $\psi$ is a harmonic spinor, and $B$ is flat, the Weitzenb\"ock formula on $Y$ implies
\[
\nabla_B^*\nabla_B \psi=\frac{1}{2}\psi,
\]
hence the pointwise identity
\[
\Delta|\psi|^2=2\langle \psi, \nabla_B^*\nabla_B \psi\rangle-2|\nabla_B\psi|^2=| \psi|^2-2|\nabla_B\psi|^2
\]
holds. Multiplying by $|\psi|^2$ and integrating, we obtain
\[
\int | \psi|^4-\int 2| \psi|^2|\nabla_B\psi|^2=\int | \psi|^2\Delta|\psi|^2\geq0.
\]
Recalling now that $|\alpha|^2=\frac{1}{4}|\psi|^4$, we obtain, by using the Bochner formula and $\lambda_1^*=1$, the chain of inequalities
\begin{align*}
2\|\alpha\|^2_{L^2}=\int\frac{1}{2}|\psi|^4\geq \int | \psi|^2|\nabla_B\psi|^2&=\|\nabla\alpha\|^2_{L^2}+\|d\alpha\|^2_{L^2}\\
&=2\|d\alpha\|^2_{L^2}-\mathrm{Ric}(\alpha,\alpha)\geq2\|d\alpha\|^2_{L^2}\geq 2\|\alpha\|^2_{L^2}.
\end{align*}
This implies that all inequalities are equalities, so that in particular $\alpha$ is a $1$-eigenform, i.e. a multiple of $\eta$. So, according to the sign of $\alpha(v)$ at $z=0$, the underlying spin$^c$ structure is either $\spin_0$ or its conjugate $\bar{\spin}_0$.
\end{proof}

We need to understand more in detail the spin$^c$ structure $\spin_0$ on $Y$; before doing this, let us study the spin geometry of the double cover $\tilde{Y}$. The manifold $\tilde{Y}=\Gamma\setminus\mathrm{Solv}$ comes with a natural spin structure $\spin_*$ coming from the left invariant orthonormal framing dual to $\Z,\X,\Y$, i.e.
\[
e_1=\frac{d}{dz},\quad e_2=e^{-z}\frac{d}{dx},\quad e_3=e^z\frac{d}{dy}.
\]
This defines a spin structure $\spin_*$ by taking the trivial bundle $S=Y\times \mathbb{C}^2$ and letting these vector fields act via the Pauli matrices
\[
\sigma_1=
\begin{bmatrix}
i&0\\
0&-i
\end{bmatrix}
\quad
\sigma_2=
\begin{bmatrix}
0&-1\\
1&0
\end{bmatrix}
\quad
\sigma_3=
\begin{bmatrix}
0&i\\
i&0
\end{bmatrix}.
\]
Let $B_*$ the spin connection on $Y$ induced by the Levi--Civita connection.



\begin{lemm}\label{harmonic}
The kernel of the Dirac operator $D_{B_*}$ consists of the constant spinors.
\end{lemm}


\begin{proof}
Let us write explicitly the Dirac operator. Our orthonormal frame satisfies the commutation relations
\begin{align*}
[e_1,e_2]&=-e_2\\
[e_1,e_3]&=e_3\\
[e_2,e_3]&=0.
\end{align*}
Setting $[e_i,e_j]=\sum_k C_{ijk} e_k$, we have that the Christoffel symbols are
\[
\Gamma_{ijk}=\frac{1}{2}(C_{ijk}-C_{ikj}-C_{jki}),
\]
hence in our case the non-zero ones are
\[
\Gamma_{212}=-\Gamma_{221}=1,\quad \Gamma_{313}=-\Gamma_{331}=-1.
\]
The spin connection on the spinor bundle is given by
\[
\nabla_{e_i}\Psi=e_i(\Psi)+\frac{1}{4}\sum_{j<k}\Gamma_{ijk}[\sigma_j,\sigma_k]\cdot \Psi,
\]
see Equation $(3.13)$ in~\cite[Section~3.3]{BGV}. Therefore, as $[\sigma_1,\sigma_2]=-2\sigma_3$ and $[\sigma_1,\sigma_3]=2\sigma_2$, we have
\begin{align*}
\nabla_{e_1}\Psi&=e_1(\Psi)\\
\nabla_{e_2}\Psi&=e_2(\Psi)-\frac{1}{2}\sigma_3\cdot\Psi\\
\nabla_{e_3}\Psi&=e_3(\Psi)-\frac{1}{2}\sigma_2\cdot\Psi
\end{align*}
As $\sigma_2$ and $\sigma_3$ anticommute, we have
\[
D_{B_*}\Psi=\sum_i\rho(e_i)\cdot\nabla_{e_i}\Psi=\sum_i\rho(e_i)\cdot{e_i}(\Psi).
\]
Hence, writing $\Psi=(f, g)$, we have
\[
D_{B_*}
\begin{bmatrix}f\\g
\end{bmatrix}=
\begin{bmatrix}
if_z-e^{-z}g_x+ie^zg_y\\
-ig_z+e^{-z}f_x+ie^zf_y
\end{bmatrix},
\]
and the equations for a harmonic spinor are
\begin{align*}
f_z+ie^{-z}g_x+e^zg_y&=0\\
g_z+ie^{-z}f_x-e^zf_y&=0.
\end{align*}
Let us now decompose the equations according to the eigenmodes $\underline{\mu}=(\mu,\mu')\in\Lambda'$. We obtain
\begin{align*}
f_z-\mu e^{-z}g+i{\mu'}e^zg&=0\\
g_z-\mu e^{-z}f-i{\mu'}e^zf&=0. \numberthis \label{ode}
\end{align*}
Of course for the zero mode the kernel consists of constant solutions. Let us show now that the eigenmodes with $\underline{\mu}\neq 0$ do not admit non-zero harmonic spinors. We~have
\begin{align*}
\frac{d}{dz}|f|^2&=2\Re\left(f_z\bar{f}\right)=2\Rel\left((\mu e^{-z}g-i\mu'e^zg)\bar{f}\right)\\
\frac{d}{dz}|g|^2&=2\Rel\left(\bar{g}_z{g}\right)=2\Rel\left(\left(\mu e^{-z}\bar{f}-i\mu'e^z\bar{f}\right)g\right)
\end{align*}
hence
\[
\frac{d}{dz}(|f|^2-|g|^2)=0.
\]
As $|f|^2-|g|^2$ is in the class of function $\mathcal{S}$ from Equation~\eqref{schwartz}, we have $|f|^2=|g|^2$ everywhere. This, together with our ODE~\eqref{ode}, shows that the functions $f$ and $g$ are never zero. We then have
\begin{align*}
f_z\bar{g}&=\mu e^{-z}|g|^2-i\mu'e^{z}|g|^2\\
f\bar{g}_z&=\mu e^{-z}|f|^2-i\mu'e^{z}|f|^2
\end{align*}
hence
\[
\frac{d}{dz}\left(\frac{f}{\bar{g}}\right)=\frac{f_z\bar{g}-f\bar{g}_z}{\bar{g}^2}=0.
\]
Therefore, up to multiplying $f$ and $g$ by the same complex constant, we have
\[
g=\bar{f}
\]
and both are equations are equivalent to
\[
f_z=\mu e^{-z}\bar{f}-i\mu'e^{z}\bar{f}.
\]
Writing $f=a+ib$ for real functions $a,b$, this can be written as the system
\begin{align*}
a_z&=\mu e^{-z}a-\mu'e^zb\\
b_z&=-\mu'e^za-\mu e^{-z}b
\end{align*}
Differentiating the first equation, and making some simple substitutions, we obtain the equation
\[
a_{zz}=a_z+\left(\mu^2e^{-2z}+\mu'^2e^{2z}-2\mu e^{-z}\right)a.
\]
Then, $A=e^{-z/2}a$ (which still lies in $\mathcal{S}$) satisfies the ODE
\[
A_{zz}=\left(\mu^2e^{-2z}+\mu'^2e^{2z}-2\mu e^{-z}+\frac{1}{4}\right)\cdot A.
\]
The claim is that under our assumption $|\mu\mu'|>8$, the first factor on the right hand side is always strictly positive, so that by Lemma~\ref{keyODE}, $A$ is zero, and so are $a$ and $b$. We need to show
\[
\mu^2e^{-2z}+\mu'^2e^{2z}-2\mu e^{-z}+\frac{1}{4}>0,
\]
which is equivalent to
\[
(\mu e^{-z}-1)^2+(\mu'e^z)^2>\frac{3}{4}.
\]
If $|\mu e^{-z}|\geq 2$, the inequality is clearly true; otherwise, we have
\[
|\mu'e^z|\geq \frac{|\mu e^{-z}|}{2}\cdot |\mu' e^z|=\frac{|\mu\mu'|}{2}\geq 4,
\]
so we are done.
\end{proof}


With this computation in mind, we will show that the Dirac operator on our rational homology sphere $Y$ equipped with the spin$^c$ structure $\spin_0$ has no kernel (the same will hold for $\bar{\spin}_0$).


First of all, we pull it back to $\tilde{Y}$; suppose that this is the mapping torus of $A\in\SL(2;\mathbb{Z})$. The Mayer--Vietoris sequence for the mapping torus of any map $f$ implies the exact sequence
\[
H_1\left(T^2;\mathbb{Z}/2\mathbb{Z}\right)\stackrel{1-f_*}{\rightarrow} H_1\left(T^2;\mathbb{Z}/2\mathbb{Z}\right)\rightarrow H_1\left(M_f;\mathbb{Z}/2\mathbb{Z}\right)\rightarrow \mathbb{Z}/2\mathbb{Z}\rightarrow 0.
\]
As $A$ is congruent to the identity modulo $2$ (see Remark~\ref{mod2}), we have that
\[
H^1\left(\tilde{Y};\mathbb{Z}/2\mathbb{Z}\right)\;\cong\;\mathbb{Z}/2\mathbb{Z}\quad\oplus\quad H^1\left(T^2;\mathbb{Z}/2\mathbb{Z}\right)\;\cong\;(\mathbb{Z}/2\mathbb{Z})^3,
\]
so that, from the point of view of spin topology, $\tilde{Y}$ looks like the more familiar three-torus. The pullback of $\spin_0$ to $\tilde{Y}$, call it $\tilde{\spin}$, also corresponds to the $1$-form $\eta$. Denoting by $\xi$ the left invariant $1$-form obtained from $\eta$ by $\pi/2$ counterclockwise rotation within the fibers, we see that the cover involution $\tau$ sends the frame $\eta,\xi,\mathcal{Z}$ to $\eta,-\xi,-\mathcal{Z}$. Therefore $\tilde{\spin}$ is the spin structure obtained from the standard one $\spin_*$ by twisting by $2\pi$ around the class dual to $[v]\in H^1(T^2;\mathbb{Z}/2\mathbb{Z})$. In particular the holonomy of the pullback of the flat connection of $\spin_0$ is trivial around the generator of $b_1(\tilde{Y})$. The sublattice of $\Lambda$ spanned by $2v$ and $w$ is preserved by $A$; the corresponding mapping torus ${\overline{Y}}$ is a double cover of $\tilde{Y}$; and the pullback of $\tilde{\spin}$ is the standard spin structure $\spin_*$ on ${\overline{Y}}$. One can then identify the harmonic spinors on $(\tilde{Y},\tilde{\spin})$ as the harmonic spinors on $({\overline{Y}},\spin_*)$ which change sign under translation by $v$ at $z=0$; by Lemma~\ref{harmonic}, there are no such spinors. Hence, there are no harmonic spinors on the base space $(Y,\spin_0)$.


Putting pieces together, we finally conclude.


\begin{proof}[Proof of Theorem~\ref{main}]By the discussion above, we have found small perturbations for which there are no irreducible solutions and the (perturbed) Dirac operator of the reducible solution has no kernel; we can then add a further small perturbation to make all of its eigenvalues simple (while preserving these properties) as in~\cite[Chapter~12]{KM}; the proof of Theorem~\ref{main} is then completed.
\end{proof}

\subsection*{Acknowledgements}
The author would like to thank Liam Watson for some helpful comments, and the anonymous referee for carefully reading the manuscript and providing very detailed feedback.

